\documentclass[9pt,a4paper]{article}
\usepackage[a4paper,margin=9mm]{geometry}
\usepackage[T1]{fontenc}
\IfFileExists{lmodern.sty}{\usepackage{lmodern}\usepackage{microtype}}{\usepackage[expansion=false]{microtype}}
\usepackage{amsmath,amssymb,booktabs,array,pdflscape}
\setlength{\parindent}{0pt}
\newcommand{\NBS}{\mathrm{NBS}}
\begin{document}\begin{landscape}\thispagestyle{empty}
{\Large\bfseries Table 2: Sensitivity of recurrent deals to the contemporaneous NBS}\\
{\small 10,000 simultaneous Monte Carlo perturbations of the stakeholder preference landscape (per-issue interval $\pm\max(1\text{ point},\,25\%$ of the issue weight$)$, total budget preserved) and of the acceptance threshold ($\pm10\%$ of $r=65$, drawn as a single shared draw applied to all stakeholders); baseline NBS: A3,B2,C3,D1,E1,F2.}\par\vspace{2mm}
\centering\scriptsize\setlength{\tabcolsep}{2.3pt}\renewcommand{\arraystretch}{1.16}
\begin{tabular}{@{}llccc|cc|cc|cc|cc@{}}
\toprule
&&&&&\multicolumn{2}{c|}{\textbf{Hamming $H$}}&\multicolumn{2}{c|}{\textbf{$L_1$}}&\multicolumn{2}{c|}{\textbf{$L_2$}}&\multicolumn{2}{c}{\textbf{Nash-efficiency loss $\mathcal L_{\NBS}$ (\%)}}\\
\cmidrule(lr){6-7}\cmidrule(lr){8-9}\cmidrule(lr){10-11}\cmidrule(lr){12-13}
\textbf{ID}&\textbf{Deal $(A,B,C,D,E,F)$}&\textbf{Runs}&\textbf{IR in MC}&\textbf{Pareto in MC}&\textbf{Base}&\textbf{MC med. [5--95\%]}&\textbf{Base}&\textbf{MC med. [5--95\%]}&\textbf{Base}&\textbf{MC med. [5--95\%]}&\textbf{Base}&\textbf{MC med. [5--95\%]}\\
\midrule
N0 & $(3,2,3,1,1,2)$ & 0 & 92.2\% & 100.0\% & 0 & 2 [0--4] & 0.0 & 22.7 [0.0--32.8] & 0.0 & 12.6 [0.0--17.5] & 0.0 & 15.2 [0.0--74.9]\\
R1 & $(3,2,3,1,2,1)$ & 5 & 34.6\% & 100.0\% & 2 & 3 [1--4] & 34.0 & 42.0 [30.8--58.9] & 20.4 & 22.0 [17.9--28.9] & -- & 76.2 [42.8--97.7]\\
R2 & $(3,2,3,1,2,2)$ & 3 & 4.7\% & 100.0\% & 1 & 3 [1--4] & 33.0 & 44.6 [30.6--62.5] & 21.5 & 24.0 [19.0--31.8] & -- & 88.1 [59.7--99.2]\\
R3 & $(3,2,3,1,2,3)$ & 3 & 55.2\% & 100.0\% & 2 & 3 [2--4] & 44.0 & 45.1 [31.8--70.6] & 26.1 & 24.3 [16.6--36.3] & 95.5 & 85.9 [58.4--99.1]\\
R4 & $(3,3,2,2,3,1)$ & 3 & 0.0\% & 100.0\% & 5 & 4 [3--5] & 95.0 & 100.7 [69.1--116.1] & 46.1 & 48.9 [36.0--54.7] & -- & --\\
R5 & $(2,2,2,1,2,2)$ & 2 & 0.0\% & 0.0\% & 3 & 4 [3--5] & 50.0 & 61.4 [47.7--79.1] & 34.3 & 37.9 [31.9--45.0] & -- & --\\
R6 & $(3,2,3,1,3,3)$ & 2 & 76.1\% & 100.0\% & 2 & 3 [2--4] & 42.0 & 41.2 [29.6--60.0] & 23.7 & 22.0 [15.3--31.7] & 81.1 & 81.0 [51.2--98.7]\\
R7 & $(4,2,3,1,2,3)$ & 2 & 0.0\% & 100.0\% & 3 & 3 [2--4] & 86.0 & 80.0 [66.6--110.6] & 46.3 & 42.9 [35.2--58.1] & -- & --\\
\bottomrule
\end{tabular}
\par\raggedright\vspace{2mm}
\textbf{Perturbation rule.} Each issue weight $w_{pi}$ is redrawn uniformly from its own admissible interval $[\max(0,w_{pi}-c_{pi}),\,w_{pi}+c_{pi}]$ with $c_{pi}=\max(1,\,0.25\,w_{pi})$, so tolerance scales with the stated weight and never falls below the 1-point granularity of the elicitation instrument. A baseline weight of $0$ is treated as ``could plausibly have been 1'' and redrawn from $[0,1]$; since all options in such an issue were scored $0$, no within-issue preference was expressed and every option receives the perturbed issue maximum, which shifts the stakeholder's utility level without altering their ranking of packages. Weights are drawn by hit-and-run, uniformly on the set of vectors that stay inside every issue interval and preserve the stakeholder's original total budget (100 points); the same sampler is used in the companion Campylobacter feasibility search. All option utilities within a non-zero issue are rescaled proportionally, so within-issue ratios, rankings, ties, and structural zeros are preserved exactly. Issue-importance reversals are therefore possible only between issues whose intervals overlap; across all draws the largest relative weight gap ever reversed was 36\%, over 20 issue pairs. The acceptance threshold is established jointly for all participants and is therefore perturbed as a single shared draw: $r'_p=r_p(1+u)$ for every stakeholder $p$, with one $u\sim U(-0.1,0.1)$ drawn per sample.
\par\vspace{1mm}
\textbf{Definitions.} For perturbation $s$ with contemporaneous NBS $d_s^*$ and thresholds $r^{(s)}$,
\[H_s(d)=\sum_j \mathbf 1[d_j\neq d^*_{s,j}],\quad L_{1,s}(d)=\sum_i|U_i^{(s)}(d)-U_i^{(s)}(d_s^*)|,\]
\[L_{2,s}(d)=\sqrt{\sum_i\left(U_i^{(s)}(d)-U_i^{(s)}(d_s^*)\right)^2},\quad
\mathcal L_{\NBS,s}(d)=100\left[1-\frac{\prod_i\left(U_i^{(s)}(d)-r_i^{(s)}\right)}{\prod_i\left(U_i^{(s)}(d_s^*)-r_i^{(s)}\right)}\right].\]
\textbf{Runs} counts how many of the 35 simulation runs put the deal forward in their final public proposal; recurrent deals are listed in descending order of that count, with the deal string as a tie-break. Every distinct deal in a run's public \texttt{<ANSWER>} block is counted, so a run proposing a primary deal and a fallback contributes both; deals appearing only in the private \texttt{<SCRATCHPAD>} are excluded. N0 is the baseline Nash bargaining solution, included because every distance is measured against it, and was not itself proposed in any run. IR and Pareto shares are over all 10,000 draws. Because a raised threshold can empty the individually rational set, the NBS may fail to exist: this occurred in 0.0\% of draws, and the $H$, $L_1$, $L_2$ and Nash-loss columns are computed over the remaining draws. Brackets are empirical Monte Carlo robustness intervals (5th--95th percentiles), not sampling-theory confidence intervals. Nash loss is summarized only when enough perturbed draws keep the target deal individually rational, and is then conditional on that subset.
\par\vspace{3mm}
\textbf{Alternative NBS distribution} (over draws in which an NBS exists).\par\vspace{1mm}
\begin{tabular}{@{}lrrr@{}}\toprule \textbf{Deal}&\textbf{Count}&\textbf{Frequency}&\textbf{$H$ vs baseline}\\\midrule
$(3,2,3,1,1,2)$ & 2499 & 25.0\% & 0\\
$(4,2,3,2,3,2)$ & 1862 & 18.6\% & 3\\
$(3,2,3,2,1,3)$ & 1614 & 16.1\% & 2\\
$(3,2,3,2,1,2)$ & 1061 & 10.6\% & 1\\
$(4,2,3,2,3,1)$ & 908 & 9.1\% & 4\\
$(3,2,3,2,1,1)$ & 800 & 8.0\% & 2\\
$(3,2,3,1,1,1)$ & 558 & 5.6\% & 1\\
$(4,2,3,2,2,1)$ & 529 & 5.3\% & 4\\
$(4,2,3,2,1,2)$ & 60 & 0.6\% & 2\\
$(3,2,3,2,2,3)$ & 37 & 0.4\% & 3\\
$(3,2,3,2,5,3)$ & 26 & 0.3\% & 3\\
$(4,2,3,2,2,2)$ & 20 & 0.2\% & 3\\
$(3,2,3,2,3,3)$ & 13 & 0.1\% & 3\\
$(4,2,3,2,5,2)$ & 10 & 0.1\% & 3\\
$(4,2,3,2,5,1)$ & 1 & 0.0\% & 4\\
$(4,2,3,2,1,1)$ & 1 & 0.0\% & 3\\
\bottomrule\end{tabular}
\end{landscape}\end{document}